\documentclass[11pt]{article}
\usepackage[a4paper,margin=2cm]{geometry}
\usepackage{booktabs}
\usepackage{array}
\usepackage{enumitem}
\usepackage{tabularx}
\usepackage{xcolor}
\setlist{nosep,leftmargin=1.35em}
\setlength{\parskip}{0.4em}
\setlength{\parindent}{0pt}
\pagestyle{empty}
\renewcommand{\arraystretch}{1.15}

\newcolumntype{Y}{>{\raggedright\arraybackslash}X}

\begin{document}

\begin{center}
    {\large\bfseries CGBuilder2: Summary of Features Added Relative to CGBuilder}
\end{center}

\vspace{-0.4em}

CGBuilder2 is an extended, modernised rewrite of the browser-based \textit{CGBuilder}
tool for constructing Martini coarse-grained (CG) molecular models. The original
CGBuilder lets a user load an atomistic structure, group atoms into CG beads by
clicking in a 3D viewer, and export index (\texttt{.ndx}), mapping (\texttt{.map}),
and coordinate (\texttt{.gro}) files. CGBuilder2 keeps this mapping workflow but
turns the tool into a complete CG \emph{topology} designer: it assembles a full
GROMACS/Martini \texttt{.itp} topology live as the model is built, and adds bonded
interactions, virtual sites, an elastic network, automatic mass handling, and
SMILES-based molecule input. It was produced as a \textbf{human-guided agentic
build}---the implementation was carried out iteratively by an AI coding agent under
close human direction---as one of several small tools developed for this thesis.

\begin{table}[h]
\centering
\footnotesize
\begin{tabularx}{\textwidth}{@{}l Y Y@{}}
\toprule
\textbf{Aspect} & \textbf{Original CGBuilder} & \textbf{CGBuilder2 (this work)} \\
\midrule
Molecule input      & File upload (PDB/GRO/CIF) & File upload \emph{plus} SMILES $\rightarrow$ 3D conformer generated in-browser \\
Bead mapping        & Click atoms into beads & Same, plus per-bead Martini type, charge, and a settable residue name \\
Topology output     & None & Live Martini \texttt{.itp} (updates continuously on the left of the display) \\
Bonds               & None & Manual bonds, built by clicking beads in 3D or by dropdown \\
Constraints         & None & Rigid constraints emitted as \texttt{\#ifndef FLEXIBLE} with a matching flexible bond under \texttt{\#ifdef FLEXIBLE} \\
Angles / dihedrals  & None & Interactive (click) and dropdown builders with editable force constants \\
Virtual sites       & None & \texttt{virtual\_sites3}: pick a 3-bead constructor triad, then click beads to virtualise \\
Elastic network     & None & Generator over bead pairs within a cutoff, with tunable strength and distance decay \\
Exclusions          & None & Generated automatically, de-duplicated and correctly formatted \\
Bead mass           & None & Automatic: sum of atom masses, mass of unmapped atoms spread along the bond graph, virtual-site mass redistributed to constructors (mass conserved) \\
Visualisation       & CG spheres, atom labels & Adds a CG-only view (bead name/number labels; virtual-site constructors in red) and click-to-select beads \\
File I/O            & Export \texttt{.ndx}, \texttt{.map}, \texttt{.gro} & Adds \texttt{.itp}, an all-atom \texttt{.gro} reference, and re-import of a CGBuilder \texttt{.ndx} mapping \\
Software base       & Single JavaScript file, CDN scripts & Modular ES modules bundled with Vite/npm \\
\bottomrule
\end{tabularx}
\end{table}

\vspace{-0.6em}

\textbf{Highlights.} The most significant additions are (i) end-to-end generation of
a simulation-ready Martini topology rather than only a mapping; (ii) mass-conserving
automatic bead masses, including redistribution of virtual-site mass onto its
constructors; (iii) two complementary ways to define every bonded term---direct
3D click-selection and dropdown menus---and (iv) in-browser SMILES input, removing
the need to pre-build an atomistic structure file.

\textbf{Implementation.} The tool runs entirely client-side. Molecular
visualisation and atom/bead picking use NGL; SMILES parsing and 3D conformer
generation use OpenChemLib; the application is written as small ES modules and
bundled with Vite. No server or network access is required at run time.

\textbf{Provenance.} CGBuilder2 was developed as a human-guided agentic build: the
author specified requirements, reviewed and steered the design decisions, and
validated the output, while an AI coding agent performed the implementation across
successive iterations. It is included here as a representative example of the
agentically-assisted tooling produced during this thesis.

\end{document}
